%! Exponential and Logarithmic Equations and Inequalities — Unit 5, Lesson 5.6 — Homework (Answer Key)
\documentclass[10pt]{article}
\usepackage{precalculus-article}
\usepackage{precalculus-key}   % pulls in -boxes; do NOT also load -boxes
\usepackage{grading}
\gradingassignment{precalc-u05-l06-hw}
\gradingcourse{precalculus}
\newcommand{\TallMath}[1]{$\displaystyle #1\rule[-1.4em]{0pt}{3.2em}$}

\begin{document}

\pageheader{Unit 5, Lesson 5.6}{Homework --- Answer Key}

\vspace{0.1in}

\begin{enumerate}[leftmargin=1.8em, itemsep=14pt]

\item Solve each equation for $x$. Where necessary, give both exact and decimal forms
      (rounded to 3 decimal places).

\begin{enumerate}[leftmargin=1.5em, itemsep=8pt, label=\alph*.]
  \item $3^x = 243$

  \vspace{4pt}
  $x = $ \smash{\gradedblank[id=hw1a, type=short, points=1, answer={5}]{4cm}{13pt}} \quad (since $3^5 = 243$)

  \item $2^x = 50$

  \vspace{4pt}
  Exact: $x = $ \smash{\gradedblank[id=hw1b1, type=short, points=1, answer={\log 50 / \log 2}]{4cm}{13pt}} \qquad Decimal: $x \approx $ \smash{\gradedblank[id=hw1b2, type=short, check=numeric, tolerance=0.0005, points=1, answer={5.644}]{2.5cm}{13pt}}\vphantom{\blank{0pt}}

  \item $7^x = \dfrac{1}{49}$

  \vspace{4pt}
  $x = $ \smash{\gradedblank[id=hw1c, type=short, points=1, answer={-2}]{4cm}{13pt}} \quad (since $7^{-2} = 1/49$)
\end{enumerate}

\item Solve each logarithmic equation. Show all steps and check for extraneous solutions.

\begin{enumerate}[leftmargin=1.5em, itemsep=10pt, label=\alph*.]
  \item $\log_2(x+3) = 5$

  \vspace{4pt}
  Work: \ansline{$x + 3 = 2^5 = 32 \Rightarrow x = 29$. Check: $\log_2(32) = 5$. Valid.}

  \vspace{4pt}
  Solution: $x = $ \smash{\gradedblank[id=hw2a1, type=short, points=1, answer={29}]{2.5cm}{13pt}} \qquad Valid? \smash{\gradedblank[id=hw2a2, type=short, check=text, points=1, answer={yes}]{2cm}{13pt}}

  \item $\log_3 x + \log_3(x-8) = 2$

  \vspace{4pt}
  Step 1 (combine): \ansline{$\log_3(x(x-8)) = 2$}

  \vspace{4pt}
  Step 2 (exponential form): \ansline{$x(x-8) = 3^2 = 9$}

  \vspace{4pt}
  Step 3 (solve quadratic): \ansline{$x^2 - 8x - 9 = 0 \Rightarrow (x-9)(x+1) = 0$}

  \vspace{4pt}
  Possible solutions: $x = $ \smash{\gradedblank[id=hw2b1, type=short, points=1, answer={9}]{2cm}{13pt}} or $x = $ \smash{\gradedblank[id=hw2b2, type=short, points=1, answer={-1}]{2cm}{13pt}}

  \vspace{4pt}
  Check for extraneous solutions: \ansline{$x=-1$: $\log_3(-1)$ undefined. Extraneous. $x=9$: $\log_3 9 + \log_3 1 = 2 + 0 = 2$. Valid.}

  \vspace{4pt}
  Final answer: $x = $ \smash{\gradedblank[id=hw2b3, type=short, points=1, answer={9}]{2.5cm}{13pt}}
\end{enumerate}

\item A population model is $P(t) = 1200 \cdot (1.04)^t$, where $t$ is in years.
      For what value of $t$ does $P(t) = 2000$? Show work using logarithms.

\vspace{4pt}
Step 1: \ansline{$1200\cdot(1.04)^t = 2000 \Rightarrow (1.04)^t = 2000/1200 = 5/3$}

\vspace{4pt}
Step 2: \ansline{$\log(1.04^t) = \log(5/3) \Rightarrow t\log 1.04 = \log(5/3)$}

\vspace{4pt}
Step 3: \ansline{$t = \log(5/3)/\log(1.04)$}

\vspace{4pt}
$t = $ \smash{\gradedblank[id=hw3a, type=short, points=2, answer={\log(5/3)/\log(1.04)}]{5cm}{13pt}} (exact) $\approx $ \smash{\gradedblank[id=hw3b, type=short, check=numeric, tolerance=0.005, points=1, answer={13.02}]{2.5cm}{13pt}} years

\item Rewrite each expression using the natural base identity $b^x = e^{(\ln b)\cdot x}$.

\begin{enumerate}[leftmargin=1.5em, itemsep=6pt, label=\alph*.]
  \item $5^x = e^{(\;\smash{\gradedblank[id=hw4a, type=short, points=1, answer={\ln 5}]{3cm}{13pt}}\;)\cdot x}$ \quad \textcolor{keyred}{$\approx e^{1.609x}$}

  \item $(0.7)^t = e^{(\;\smash{\gradedblank[id=hw4b, type=short, points=1, answer={\ln 0.7}]{3cm}{13pt}}\;)\cdot t}$ \quad \textcolor{keyred}{$\approx e^{-0.357t}$}
\end{enumerate}

\item Find the inverse of $f(x) = 2 \cdot 5^{x+3} - 1$. Show each step. Then state the
      domain and range of both $f$ and $f^{-1}$.

\vspace{6pt}
\ansline{Swap: $x = 2\cdot5^{y+3}-1$. Add 1: $x+1=2\cdot5^{y+3}$. Divide by 2: $\tfrac{x+1}{2}=5^{y+3}$.\\[3pt]
Apply $\log_5$: $\log_5\tfrac{x+1}{2}=y+3$. Subtract 3: $y=\log_5\tfrac{x+1}{2}-3$.}

\vspace{4pt}
$f^{-1}(x) = $ \smash{\gradedblank[id=hw5a, type=short, points=2, answer={\log_5\frac{x+1}{2} - 3}, display={$\log_5((x+1)/2) - 3$}]{6cm}{13pt}}

\vspace{6pt}
Domain of $f$: \smash{\gradedblank[id=hw5b, type=short, check=text, points=1, answer={all reals}]{3.5cm}{13pt}} \qquad Range of $f$: \smash{\gradedblank[id=hw5c, type=short, check=text, points=1, answer={(-1, \infty)}, display={$(-1,\infty)$}]{3.5cm}{13pt}}

\vspace{4pt}
Domain of $f^{-1}$: \smash{\gradedblank[id=hw5d, type=short, check=text, points=1, answer={(-1, \infty)}, display={$(-1,\infty)$}]{3.5cm}{13pt}} \qquad Range of $f^{-1}$: \smash{\gradedblank[id=hw5e, type=short, check=text, points=1, answer={all reals}]{3.5cm}{13pt}}

\item \textbf{AP-style multiple choice.} Which value of $x$ satisfies
      $\log_2(x+4) + \log_2 x = 5$?

\vspace{6pt}
\begin{enumerate}[leftmargin=2em, itemsep=4pt, label=(\Alph*)]
  \item $x = -8$
  \item $x = 4$ \quad \textcolor{keyred}{\textbf{$\leftarrow$ correct}}
  \item $x = 2$
  \item $x = 8$ \hfill\textit{Letter:}~\smash{\gradedblank[id=hw6, type=choice, points=1, answer=B]{0.42in}{13pt}}
\end{enumerate}

\begin{teachernote}
Combining: $\log_2(x(x+4))=5 \Rightarrow x(x+4)=32 \Rightarrow x^2+4x-32=0
\Rightarrow (x+8)(x-4)=0$. So $x=4$ or $x=-8$.
Check: $x=-8$ gives $\log_2(-8)$ undefined. Extraneous. Answer: $x=4$.
\end{teachernote}

\end{enumerate}

\vspace{0.1in}

\begin{extensionbox}
\textbf{Extension:} Solve the inequality $\log_3(x+1) \geq 2$. Express the solution as an
inequality and justify why $x = -1$ is not in the solution set.

\ansline{$\log_3(x+1) \geq 2 \Rightarrow x+1 \geq 3^2 = 9 \Rightarrow x \geq 8$.\\[3pt]
$x=-1$ is excluded because $\log_3(-1+1)=\log_3 0$ is undefined (domain of log requires strictly positive argument).}
\end{extensionbox}

\vspace{0.1in}

\begin{reflectionbox}
What is one technique from today's lesson you feel confident about, and one situation
(equation type) where you still want more practice?

\writelines{2}
\end{reflectionbox}

\end{document}
